\chapter[Considerações Finais]{Considerações Finais}
\label{cap:consideracoes}

Este trabalho propõe uma arquitetura de MLOps para plataformas de dados
governamentais, instanciada na previsão de paralisação de obras do Programa
Minha Casa, Minha Vida no Gov Hub Cidades. A proposta parte de uma constatação
da literatura: a maior parte do esforço e da dívida técnica de um sistema de
ML está fora do código de treino \cite{sculley2015}, e no setor público esse
esforço é ainda maior por obstáculos institucionais \cite{ferreira2025}.

A principal decisão da arquitetura é estender o \textit{Data Lakehouse}
existente em vez de construir uma plataforma de ML paralela. Painel e rótulo
ficam na camada prata; features, escores e monitoramento, na camada ouro; e os
escores voltam pela camada bronze como qualquer outra fonte, sob os mesmos
testes e o mesmo catálogo das demais tabelas. Todo cálculo sobre dados fica na
camada de transformação, o ambiente de modelagem apenas treina, avalia e
escora, e o orquestrador coordena as etapas. Com isso, o treino e o escore
usam a mesma implementação de cada feature, e qualquer resultado pode ser
refeito a partir da versão do código e da versão do dataset. Como a
arquitetura é definida por papéis, e não por ferramentas, ela pode ser
reaproveitada em plataformas com outras pilhas tecnológicas.

No TCC1 foram concluídas a identificação do problema, a definição dos
requisitos, o projeto da arquitetura e o levantamento inicial dos dados. O
levantamento já revelou um achado relevante para o caso: as trocas de fonte de
dados ao longo da série produzem transições de situação que não correspondem
necessariamente a paralisações reais, o que restringe o período em que o
rótulo é confiável.

\section{Próximos passos}

Para o TCC2, estão previstas as seguintes atividades:

\begin{enumerate}
	\item[A1] Concluir a análise exploratória e fechar as decisões sobre
	horizonte, capacidade, retomada e usuário do alerta;
	\item[A2] Calcular o desempenho da regra vigente no período SFTP;
	\item[A3] Implementar na camada de transformação o painel e o rótulo
	(camada prata) e as features \textit{as of} $t$ (camada ouro), com testes
	de vazamento;
	\item[A4] Implementar o export de datasets versionados, o treino e a
	avaliação, com registro dos experimentos;
	\item[A5] Implementar a promoção com \textit{gate}, o escore mensal e o
	monitoramento;
	\item[A6] Executar a avaliação (AV1--AV5);
	\item[A7] Redigir a monografia final e preparar a defesa.
\end{enumerate}

\section{Cronograma}

A Tabela~\ref{tab:cronograma} apresenta o cronograma previsto para o TCC2.

\begin{table}[htb]
	\centering
	\caption{Cronograma das atividades do TCC2.}
	\label{tab:cronograma}
	\small
	\begin{tabular}{lcccccc}
		\toprule
		\textbf{Atividade} & \textbf{Mês 1} & \textbf{Mês 2} & \textbf{Mês 3} & \textbf{Mês 4} & \textbf{Mês 5} & \textbf{Mês 6} \\
		\midrule
		A1 -- EDA e decisões       & X & X &   &   &   &   \\
		A2 -- Regra vigente        & X &   &   &   &   &   \\
		A3 -- Painel, rótulo e features &   & X & X &   &   &   \\
		A4 -- Treino e registro    &   &   & X & X &   &   \\
		A5 -- Gate, escore e monitoramento &   &   &   & X & X &   \\
		A6 -- Avaliação            &   &   &   &   & X & X \\
		A7 -- Escrita e defesa     &   &   &   & X & X & X \\
		\bottomrule
	\end{tabular}
	\fonte{Elaborado pelos autores.}
\end{table}

\pendente{Substituir ``Mês 1'' a ``Mês 6'' pelos meses reais do semestre do
TCC2, conforme o calendário acadêmico da UnB.}
